\ifdefined\gkver\else\def\gkver{student}\fi
\documentclass[\gkver]{gaokaozhenti}
\begin{document}

\chapter{2013年安徽理综}
%% paperType: 理综物理部分

\begin{enumerate}
\item
%% number: 14
%% paperName: 2013年安徽理综第 14 题 6 分
%% typeId: 1
%% type: 单选题
%% chapter: 运动和力
%% point: 牛顿第二定律
%% method: 无
%% score: 6
%% degree: 650
%% duplicateId: 0
%% body:
如图所示，细线的一端系一质量为 $m$ 的小球，另一端固定在倾角为 $\theta$ 的光滑斜面体顶端，细线与斜面平行。在斜面体以加速度 $a$ 水平向右做匀加速直线运动的过程中，小球始终静止在斜面上，小球受到细线的拉力 $T$ 和斜面的支持力 $F_{N}$ 分别为（重力加速度为 $g$）\xzanswer{A}

\onepicture[w=6cm, align=right]{figs/14.png}

\fourchoices[answer=A]
{$T=m(g\sin\theta+a\cos\theta)$，$F_{N}=m(g\cos\theta-a\sin\theta)$}
{$T=m(g\cos\theta+a\sin\theta)$，$F_{N}=m(g\sin\theta-a\cos\theta)$}
{$T=m(a\cos\theta-g\sin\theta)$，$F_{N}=m(g\cos\theta+a\sin\theta)$}
{$T=m(a\sin\theta-g\cos\theta)$，$F_{N}=m(g\sin\theta+a\cos\theta)$}

%% answer: A
%% memo:
\memoanswer{
小球受力如图，根据牛顿第二定律：沿斜面方向有 $T-mg\sin\theta=ma\cos\theta$；垂直于斜面方向上有 $mg\cos\theta-F_{N}=ma\sin\theta$；解得 $T=m(g\sin\theta+a\cos\theta)$，$F_{N}=m(g\cos\theta-a\sin\theta)$。

\onepicture[w=5cm]{figs/14_an.png}

正确选项：A。
}

\item
%% number: 15
%% paperName: 2013年安徽理综第 15 题 6 分
%% typeId: 1
%% type: 单选题
%% chapter: 磁场
%% point: 洛伦兹力
%% method: 无
%% score: 6
%% degree: 850
%% duplicateId: 0
%% body:
图中 a、b、c、d 为四根与纸面垂直的长直导线，其横截面积位于正方形的四个顶点上，导线中通有大小相等的电流，方向如图所示。一带正电的粒子从正方形中心 O 点沿垂直于纸面的方向向外运动，它所受洛伦兹力的方向是\xzanswer{B}

\onepicture[w=5cm, align=right]{figs/15.png}

\fourchoices[answer=B]
{向上}
{向下}
{向左}
{向右}

%% answer: B
%% memo:
\memoanswer{
根据右手安培定则可判定 O 点磁感应强度的方向水平向左，根据左手定则可判定：一带正电的粒子从正方形中心 O 点沿垂直于纸面的方向向外运动，它所受洛伦兹力的方向是向下。

正确选项：B
}

\item
%% number: 16
%% paperName: 2013年安徽理综第 16 题 6 分
%% typeId: 1
%% type: 单选题
%% chapter: 电磁感应
%% point: 电磁感应中的变加速
%% method: 无
%% score: 6
%% degree: 650
%% duplicateId: 0
%% body:
如图所示，足够长平行金属导轨倾斜放置，倾角为 $37^{\circ}$，宽度为 $0.5\Um$，电阻忽略不计，其上端接一小灯泡，电阻为 $1\UO$。一导体棒 MN 垂直于导轨放置，质量为 $0.2\Ukg$，接入电路的电阻为 $1\UO$，两端与导轨接触良好，与导轨间的动摩擦因数为 0.5。在导轨间存在着垂直于导轨平面的匀强磁场，磁感应强度为 $0.8\UT$。将导体棒 MN 由静止释放，运动一段时间后，小灯泡稳定发光，此后导体棒 MN 的运动速度及小灯泡消耗的电功率分别为（重力加速度 $g$ 取 $10\Umsq$，$\sin37^{\circ}=0.6$）\xzanswer{B}

\onepicture[w=6cm, align=right]{figs/16.png}

\fourchoices[answer=B]
{$2.5\Ums$，$1\UW$}
{$5\Ums$，$1\UW$}
{$7.5\Ums$，$9\UW$}
{$15\Ums$，$9\UW$}

%% answer: B
%% memo:
\memoanswer{
小灯泡稳定发光，说明导体棒 MN 匀速运动，则有 $mg\sin37^{\circ}-\frac{B^2L^2v}{R_L+R}-\mu mg\cos\theta=0$，可得导体棒 MN 的运动速度 $v=5\Ums$；小灯泡消耗的电功率 $P=I^{2}R_{L}=(\frac{BLv}{R_L+R})^2R_L=1\UW$。

正确选项：B。
}

\item
%% number: 17
%% paperName: 2013年安徽理综第 17 题 6 分
%% typeId: 1
%% type: 单选题
%% chapter: 曲线运动
%% point: 人造地球卫星
%% method: 无
%% score: 6
%% degree: 650
%% duplicateId: 0
%% body:
质量为 $m$ 的人造地球卫星与地心的距离为 $r$ 时，引力势能可表示为 $E_{p}=-\frac{GMm}{r}$，其中 $G$ 为引力常量，$M$ 为地球质量。该卫星原来在半径为 $R_{1}$ 的轨道上绕地球做匀速圆周运动，由于受到极稀薄空气的摩擦作用，飞行一段时间后其圆周运动的半径变为 $R_{2}$，此过程中因摩擦而产生的热量为\xzanswer{C}

\fourchoices[answer=C]
{$GMm(\frac{1}{R_2}-\frac{1}{R_1})$}
{$GMm(\frac{1}{R_1}-\frac{1}{R_2})$}
{$\frac{GMm}{2}(\frac{1}{R_2}-\frac{1}{R_1})$}
{$\frac{GMm}{2}(\frac{1}{R_1}-\frac{1}{R_2})$}

%% answer: C
%% memo:
\memoanswer{
人造地球卫星绕地球做匀速圆周运动，万有引力提供向心力：$G\frac{Mm}{r^2}=m\frac{v^2}{r}$，故人造地球卫星的动能 $E_{k}=\frac{1}{2}mv^{2}=\frac{GMm}{2r}$，而引力势能 $E_{p}=-\frac{GMm}{r}$，人造地球卫星机械能 $E=E_{p}+E_{k}$。由能量守恒定律，因摩擦而产生的热量 $Q=E_{1}-E_{2}=(E_{k1}+E_{p1})-(E_{k2}+E_{p2})$，代入 $R_{1}$ 和 $R_{2}$ 得 $Q=\frac{GMm}{2}(\frac{1}{R_2}-\frac{1}{R_1})$。

正确选项：C。
}

\item
%% number: 18
%% paperName: 2013年安徽理综第 18 题 6 分
%% typeId: 1
%% type: 单选题
%% chapter: 曲线运动
%% point: 斜抛运动
%% method: 无
%% score: 6
%% degree: 650
%% duplicateId: 0
%% body:
由消防水龙带的喷嘴喷出水的流量是 $0.28\Umc/\Umin$，水离开喷口时的速度大小为 $16\sqrt{3}\Ums$，方向与水平面夹角为 $60^{\circ}$，在最高处正好到达着火位置，忽略空气阻力，则空中水柱的高度和水量分别是（重力加速度 $g$ 取 $10\Umsq$）\xzanswer{A}

\fourchoices[answer=A]
{$28.8\Um$，$1.12\times10^{-2}\Umc$}
{$28.8\Um$，$0.672\Umc$}
{$38.4\Um$，$1.29\times10^{-2}\Umc$}
{$38.4\Um$，$0.776\Umc$}

%% answer: A
%% memo:
\memoanswer{
水离开喷口时竖直分速度为 $v_{\perp}=v\sin60^{\circ}=24\Ums$，在竖直方向上升的高度 $H=\frac{v_{\perp}^2}{2g}=28.8\Um$，水离开喷口到达着火位置所用时间为 $t=\frac{v_{\perp}}{g}=2.4\Us$，空中水柱的水量为 $V=Qt=0.28\times\frac{2.4}{60}=1.12\times10^{-2}\Umc$。

正确选项：A。
}

\item
%% number: 19
%% paperName: 2013年安徽理综第 19 题 6 分
%% typeId: 1
%% type: 单选题
%% chapter: 电路
%% point: 电路实验
%% method: 无
%% score: 6
%% degree: 650
%% duplicateId: 0
%% body:
用图示的电路可以测量电阻的阻值。图中 $R_{x}$ 是待测电阻，$R_{0}$ 是定值电阻，G 是灵敏度很高的电流表，MN 是一段均匀的电阻丝。闭合开关，改变滑动头 P 的位置，当通过电流表 G 的电流为零时，测得 $MP=l_{1}$，$PN=l_{2}$，则 $R_{x}$ 的阻值为\xzanswer{C}

\onepicture[w=5.5cm, align=right]{figs/19.png}

\fourchoices[answer=C]
{$\frac{l_1}{l_2}R_{0}$}
{$\frac{l_1}{l_1+l_2}R_{0}$}
{$\frac{l_2}{l_1}R_{0}$}
{$\frac{l_2}{l_1+l_2}R_{0}$}

%% answer: C
%% memo:
\memoanswer{
通过电流表 G 的电流为零时，P 点的电势与 $R_{0}$ 和 $R_{x}$ 连接点的电势相等。则 $\frac{R_0}{R_x}=\frac{R_{l1}}{R_{l2}}$。又电阻丝的电阻 $R\propto l$，故 $\frac{R_0}{R_x}=\frac{l_1}{l_2}$，所以 $R_{x}=\frac{l_2}{l_1}R_{0}$。

正确选项：C。
}

\item
%% number: 20
%% paperName: 2013年安徽理综第 20 题 6 分
%% typeId: 1
%% type: 单选题
%% chapter: 电场
%% point: 静电感应
%% method: 无
%% score: 6
%% degree: 450
%% duplicateId: 0
%% body:
如图所示，$xOy$ 平面是无穷大导体的表面，该导体充满 $z<0$ 的空间，$z>0$ 的空间为真空。将电荷为 $q$ 的点电荷置于 $z$ 轴上 $z=h$ 处，则在 $xOy$ 平面上会产生感应电荷。空间任意一点处的电场皆是由点电荷 $q$ 和导体表面上的感应电荷共同激发的。已知静电平衡时导体内部场强处处为零，则在 $z$ 轴上 $z=\frac{h}{2}$ 处的场强大小为（$k$ 为静电力常量）\xzanswer{D}

\onepicture[w=5.5cm, align=right]{figs/20.png}

\fourchoices[answer=D]
{$k\frac{4q}{h^2}$}
{$k\frac{4q}{9h^2}$}
{$k\frac{32q}{9h^2}$}
{$k\frac{40q}{9h^2}$}

%% answer: D
%% memo:
\memoanswer{
由点电荷产生的电场在导体表面间的电场分布，导体表面上的感应电荷等效于在 $z=-h$ 处的带等量的异种电荷 $-q$，故在 $z$ 轴上 $z=\frac{h}{2}$ 处的场强大小为 $E=k\frac{q}{(\frac{h}{2})^2}+k\frac{q}{(\frac{h}{2}+h)^2}=k\frac{40q}{9h^2}$。

正确选项：D。
}

\item
%% number: 21
%% paperName: 2013年安徽理综第 21 题 7 分
%% typeId: 4
%% type: 实验题
%% chapter: 电路
%% point: 测电动势与内阻
%% method: 无
%% score: 7
%% degree: 650
%% duplicateId: 0
%% body:
\begin{enumerate}
	\item
	Ⅰ．（5 分）根据单摆周期公式 $T=2\pi\sqrt{\frac{l}{g}}$，可以通过实验测量当地的重力加速度。如图~\ref{2013安徽理综21a} 所示，将细线的上端固定在铁架台上，下端系一小钢球，就做成了单摆。

	（1）用游标卡尺测量小钢球直径，读数如图~\ref{2013安徽理综21b} 所示，读数为\tkanswer[2cm]{18.6}$\Umm$。

	（2）以下是实验过程中的一些做法，其中正确的有\tkanswer[2cm]{abe}。

	a．摆线要选择细些的、伸缩性小些的，并且尽可能长一些；

	b．摆球尽量选择质量大些、体积小些的；

	c．为了使摆的周期大一些，以方便测量，开始时拉开摆球，使摆线相距平衡位置有较大的角度；

	d．拉开摆球，使摆线偏离平衡位置大于 $5^{\circ}$，在释放摆球的同时开始计时，当摆球回到开始位置时停止计时，此时间间隔 $\Delta t$ 即为单摆周期 $T$；

	e．拉开摆球，使摆线偏离平衡位置不大于 $5^{\circ}$，释放摆球，当摆球振动稳定后，从平衡位置开始计时，记下摆球做 50 次全振动所用的时间 $\Delta t$，则单摆周期 $T=\frac{\Delta t}{50}$。

	\twopicture[labelA=2013安徽理综21a, labelB=2013安徽理综21b, numA=a, numB=b, wA=4.5cm, wB=7cm, align=right]
	{figs/21a.png}
	{figs/21b.png}

	\item
	Ⅱ．（6 分）（1）在测定一根粗细均匀合金丝电阻率的实验中，利用螺旋测微器测定合金丝直径的过程如图所示，校零时的读数为\tkanswer[2cm]{0.007}$\Umm$，合金丝的直径为\tkanswer[2cm]{0.638}$\Umm$。

	（2）为了精确测量合金丝的电阻 $R_{x}$，设计出如图~\ref{2013安徽理综21d} 所示的实验电路图，按照该电路图完成图~\ref{2013安徽理综21e} 中的实物电路连接。

	\onepicture[num=c, w=5.5cm, align=right]{figs/21c.png}

	\twopicture[labelA=2013安徽理综21d, labelB=2013安徽理综21e, numA=d, numB=e, wA=5cm, wB=7cm, align=right]
	{figs/21d.png}
	{figs/21e.png}

	\item
	Ⅲ．（7 分）根据闭合电路欧姆定律，用图~\ref{2013安徽理综21f} 所示电路可以测定电池的电动势和内电阻。图中 $R_{0}$ 两端的对应电压 $U_{12}$，对所得的实验数据进行处理，就可以实现测量目的。根据实验数据在 $\frac{1}{U_{12}}-R$ 坐标系中描出坐标点，如图~\ref{2013安徽理综21g} 所示。已知 $R_{0}=150\UO$，请完成以下数据分析和处理。

	（1）图~\ref{2013安徽理综21g} 中电阻为\tkanswer[2cm]{80.00}$\UO$ 的数据点应剔除；

	（2）在坐标纸上画出 $\frac{1}{U_{12}}-R$ 关系图线；

	（3）图线的斜率是\tkanswer[2.5cm]{0.00444}（$\UV^{-1}\cdot\UO^{-1}$），由此可得电池电动势 $E_{n}=$\tkanswer[2cm]{1.50}$\UV$。

	\twopicture[labelA=2013安徽理综21f, labelB=2013安徽理综21g, numA=f, numB=g, wA=4.5cm, wB=5.5cm, align=right]
	{figs/21f.png}
	{figs/21g.png}
\end{enumerate}

%% answer:
\jdanswer{
\begin{enumerate}
	\item
	Ⅰ．（1）$18.6\Umm$；（2）abe。

	\item
	Ⅱ．（1）$0.007\Umm$，$0.638\Umm$；（2）见解析（实物电路连接如图所示）。

	\item
	Ⅲ．（1）80.00；（2）见解析；（3）0.00444，$1.50\UV$。

\end{enumerate}
}
%% memo:
\memoanswer{
Ⅰ．（1）由题图主尺读数 $18\Umm$，游标尺的第 6 条刻度线与主尺刻度线对齐读数应为 $0.6\Umm$，小钢球直径 $D=18\Umm+0.6\Umm=18.6\Umm$。

（2）根据实验要求，摆线要选择细些的、伸缩性小些的，并且尽可能长一些，a 正确；为了减小空气阻力的影响，应尽量选择质量大些、体积小些的摆球，b 正确；为了使单摆的运动为简谐运动，要求摆角不大于 $5^{\circ}$，为了减小测量的误差，应使摆球振动稳定后，从平衡位置开始计时，e 正确。

Ⅱ．（1）由图可读得 $0.007\Umm$ 和 $0.645\Umm$，故校零时的读数为 $0.007\Umm$；合金丝的直径为 $D=0.645\Umm-0.007\Umm=0.638\Umm$。

（2）按照电路图连接电路如图。

\onepicture[w=6cm]{figs/21_an1.png}

Ⅲ．（1）由图 2 发现在拟合的过程中，80.00 的数据点误差太大，故应剔除。

（2）拟合作图如图所示。

\onepicture[w=5.5cm]{figs/21_an2.png}

（3）用较远的两点得图线斜率 $k=\frac{1.50-0.70}{180-0}=0.00444$，由闭合电路欧姆定律知 $E_{n}=U_{12}+\frac{U_{12}}{R_0}(R+r)$ 可得 $\frac{1}{U_{12}}=\frac{1}{E_n}+\frac{r}{E_nR_0}+\frac{1}{E_nR_0}R$，所以 $k=\frac{1}{E_nR_0}$，代入 $k$、$R_0$ 值得 $E_{n}=1.50\UV$。
}

\item
%% number: 22
%% paperName: 2013年安徽理综第 22 题 14 分
%% typeId: 6
%% type: 计算题
%% chapter: 运动和力
%% point: 牛顿定律的应用
%% method: 无
%% score: 14
%% degree: 650
%% duplicateId: 0
%% body:
一物体放在水平地面上，如图~\ref{2013安徽理综22a} 所示，已知物体所受水平拉力 $F$ 随时间 $t$ 的变化情况如图~\ref{2013安徽理综22b} 所示，物体相应的速度 $v$ 随时间 $t$ 的变化关系如图~\ref{2013安徽理综22c} 所示。求：
\begin{enumerate}
	\item
	$0\sim8\Us$ 时间内拉力的冲量；
	\item
	$0\sim6\Us$ 时间内物体的位移；
	\item
	$0\sim10\Us$ 时间内，物体克服摩擦力所做的功。
\end{enumerate}

\threepicture[labelA=2013安徽理综22a, labelB=2013安徽理综22b, labelC=2013安徽理综22c, numA=a, numB=b, numC=c, wA=5cm, wB=4.5cm, wC=4.5cm, align=right]
{figs/22a.png}
{figs/22b.png}
{figs/22c.png}

%% answer:
\jdanswer{
\begin{enumerate}
	\item
	$18\UNs$

	\item
	$6\Um$

	\item
	$30\UJ$

\end{enumerate}
}
%% memo:
\memoanswer{
（1）由图 2 知 $I=F_{1}\Delta t_{1}+F_{2}\Delta t_{2}+F_{3}\Delta t_{3}$，$I=18\UNs$。

（2）由图 3 知物体的位移为 $x=\frac{1}{2}\times4\times3\Um=6\Um$。

（3）由图 2 知，在 $6\sim8\Us$ 时间内，物体作匀速运动，于是有 $f=2\UN$；由图 3 知在 $0\sim10\Us$ 时间内物体的总位移为 $l=\frac{(8-6)+(10-2)}{2}\times3\Um=15\Um$，所以 $W=fl=2\times15\UJ=30\UJ$。
}

\item
%% number: 23
%% paperName: 2013年安徽理综第 23 题 16 分
%% typeId: 6
%% type: 计算题
%% chapter: 磁场
%% point: 带电粒子在磁场中的运动
%% method: 无
%% score: 16
%% degree: 650
%% duplicateId: 0
%% body:
如图所示的平面直角坐标系 $xOy$，在第 Ⅰ 象限内有平行于 $y$ 轴的匀强电场，方向沿 $y$ 正方向；在第 Ⅳ 象限的正三角形 abc 区域内有匀强电场，方向垂直于 $xOy$ 平面向里，正三角形边长为 $L$，且 ab 边与 $y$ 轴平行。一质量为 $m$、电荷量为 $q$ 的粒子，从 $y$ 轴上的 P（0，$h$）点，以大小为 $v_{0}$ 的速度沿 $x$ 轴正方向射入电场，通过电场后从 $x$ 轴上的 a（2$h$，0）点进入第 Ⅳ 象限，又经过磁场从 $y$ 轴上的某点进入第 Ⅲ 象限，且速度与 $y$ 轴负方向成 $45^{\circ}$ 角，不计粒子所受的重力。求：
\begin{enumerate}
	\item
	电场强度 $E$ 的大小；
	\item
	粒子到达 a 点时速度的大小和方向；
	\item
	abc 区域内磁场的磁感应强度 $B$ 的最小值。
\end{enumerate}

\onepicture[w=5.5cm, align=right]{figs/23.png}

%% answer:
\jdanswer{
\begin{enumerate}
	\item
	$\frac{mv_0^2}{2qh}$

	\item
	$\sqrt{2}v_0$，方向指向第 Ⅳ 象限与 $x$ 轴正方向成 $45^{\circ}$ 角

	\item
	$\frac{2mv_0}{qL}$

\end{enumerate}
}
%% memo:
\memoanswer{
（1）设粒子在电场中运动的时间为 $t$，则有 $x=v_{0}t=2h$，$y=\frac{1}{2}at^{2}=h$，$qE=ma$；联立以上各式可得 $E=\frac{mv_0^2}{2qh}$。

（2）粒子到达 a 点时沿负 $y$ 方向的分速度为 $v_{y}=at=v_{0}$，所以 $v=\sqrt{v_0^2+v_y^2}=\sqrt{2}v_{0}$，方向指向第 Ⅳ 象限与 $x$ 轴正方向成 $45^{\circ}$ 角。

（3）粒子在磁场中运动时，有 $qvB=m\frac{v^{2}}{r}$；当粒子从 b 点射出时，磁场的磁感应强度为最小值，此时有 $r=\frac{\sqrt{2}}{2}L$，所以 $B=\frac{2mv_0}{qL}$。

\onepicture[w=5cm]{figs/23_an.png}
}

\item
%% number: 24
%% paperName: 2013年安徽理综第 24 题 20 分
%% typeId: 6
%% type: 计算题
%% chapter: 机械振动
%% point: 简谐振动
%% method: 无
%% score: 20
%% degree: 450
%% duplicateId: 0
%% body:
如图所示，质量为 $M$、倾角为 $\alpha$ 的斜面体（斜面光滑且足够长）放在粗糙的水平地面上，底部与地面的动摩擦因数为 $\mu$，斜面顶端与劲度系数为 $k$、自然长度为 $l$ 的轻质弹簧相连，弹簧的另一端连接着质量为 $m$ 的物块。压缩弹簧使其长度为 $\frac{3}{4}l$ 时将物块由静止开始释放，且物块在以后的运动中，斜面体始终处于静止状态。重力加速度为 $g$。
\begin{enumerate}
	\item
	求物块处于平衡位置时弹簧的长度；
	\item
	选物块的平衡位置为坐标原点，沿斜面向下为正方向建立坐标轴，用 $x$ 表示物块相对于平衡位置的位移，证明物块做简谐运动；
	\item
	求弹簧的最大伸长量；
	\item
	为使斜面始终处于静止状态，动摩擦因数 $\mu$ 应满足什么条件（假设滑动摩擦力等于最大静摩擦力）？
\end{enumerate}

\onepicture[w=6cm, align=right]{figs/24.png}

%% answer:
\jdanswer{
\begin{enumerate}
	\item
	$l+\frac{mg\sin\alpha}{k}$

	\item
	当物块的位移为 $x$ 时，弹簧伸长量为 $x+\Delta l$，物块所受合力为 $F_{\text{合}}=mg\sin\alpha-k(x+\Delta l)$，联立以上各式可得 $F_{\text{合}}=-kx$，可知物块作简谐运动。

	\item
	$\frac{l}{4}+\frac{mg\sin\alpha}{k}$

	\item
	$\mu\geqslant\frac{(kl+4mg\sin\alpha)\cos\alpha}{4Mg+4mg\cos^2\alpha-kl\sin\alpha}$

\end{enumerate}
}
%% memo:
\memoanswer{
（1）设物块在斜面上平衡时，弹簧的伸长量为 $\Delta l$，有 $mg\sin\alpha-k\Delta l=0$，解得 $\Delta l=\frac{mg\sin\alpha}{k}$，此时弹簧的长度为 $l+\frac{mg\sin\alpha}{k}$。

（2）当物块的位移为 $x$ 时，弹簧伸长量为 $x+\Delta l$，物块所受合力为 $F_{\text{合}}=mg\sin\alpha-k(x+\Delta l)$，联立以上各式可得 $F_{\text{合}}=-kx$，可知物块作简谐运动。

（3）物块作简谐运动的振幅为 $A=\frac{l}{4}+\frac{mg\sin\alpha}{k}$，由对称性可知，最大伸长量为 $\frac{l}{4}+\frac{mg\sin\alpha}{k}$。

（4）设物块位移 $x$ 为正，则斜面体受力情况如图所示，由于斜面体平衡，所以有：水平方向 $f+F_{N1}\sin\alpha-F\cos\alpha=0$；竖直方向 $F_{N2}-Mg-F_{N1}\cos\alpha-F\sin\alpha=0$；又 $F=k(x+\Delta l)$，$F_{N1}=mg\cos\alpha$；联立可得 $f=kx\cos\alpha$，$F_{N2}=Mg+mg+kx\sin\alpha$。为使斜面体始终处于静止，结合牛顿第三定律，应有 $|f|\leqslant\mu F_{N2}$，所以 $\mu\geqslant\frac{|f|}{F_{N2}}=\frac{k|x|\cos\alpha}{Mg+mg+kx\sin\alpha}$；当 $x=-A$ 时，上式右端达到最大值，于是有 $\mu\geqslant\frac{(kl+4mg\sin\alpha)\cos\alpha}{4Mg+4mg\cos^2\alpha-kl\sin\alpha}$。

\onepicture[w=5cm]{figs/24_an2.png}
}

\end{enumerate}

\end{document}
